%% kapa: exam=2 round=32 year=2021 period=3 subject=compensation_law subject_ko=감정평가 및 보상법규 seq=3 number=3 points=20
%% source: Q-Net 감정평가사 기출문제 (cst00309) artlSeq=5208993 file=raw/1008/5208993/감평 2차 시험지(2021).pdf
%% conversion: kordoc 4.24.51 -> tools/gen2.py (auto); review=proofread-sonnet; flags=none
\begin{longquestion}{【문제 3】}
감정평가사 甲과 乙은 「감정평가 및 감정평가사에 관한 법률」에 따른 감정평가준칙을 위반하여 감정평가를 하였음을 이유로 업무정지처분을 받게 되었으나, 국토교통부장관은 그 업무정지처분이 「부동산 가격공시에 관한 법률」에 따른 표준지공시지가 공시 등의 업무를 정상적으로 수행하는 데에 지장을 초래할 우려가 있음을 들어, 2021. 4. 1. 甲과 乙에게 업무정지처분을 갈음하여 각 3천만원의 과징금을 부과하였다. 다음 물음에 답하시오. (20점) \par
\medskip\noindent 물음 1) 甲은 부과된 과징금이 지나치게 과중하다는 이유로 국토교통부장관에게 이의신청을 하였고, 이에 대해서 국토교통부장관은 2021. 4. 30. 甲에 대하여 과징금을 2천만원으로 감액하는 결정을 하였다. 甲은 감액된 2천만원의 과징금도 과중하다고 생각하여 과징금부과처분의 취소를 구하는 소를 제기하고자 한다. 이 경우 甲이 취소를 구하여야 하는 대상은 무엇인지 검토하시오. (10점) \par
\medskip\noindent 물음 2) 乙은 2021. 6. 1. 자신에 대한 3천만원의 과징금부과처분의 취소를 구하는 소를 제기하였다. 이에 대한 심리 결과 법원이 적정한 과징금의 액수는 1천 5백만원이라고 판단하였을 때, 법원이 내릴 수 있는 판결의 내용에 관하여 검토하시오. (10점) \par
\end{longquestion}
